\begin{exercise}{35}
Let $\mathcal C$ be a category with finite products and finite coproducts. We say that
$\mathcal C$ is \emph{pointed} if the canonical map $\varnothing\to\ast$ from the initial
to the terminal object is an isomorphism.
For the first assertion, assume that $\mathcal C$ is pointed.\footnote{The original
proceeds to ``Show that the identity canonically refines''
without this hypothesis. Insert the pointedness assumption for this assertion:
in $\mathrm{Set}$ the lax structure at $(\ast,\varnothing)$
would require a function
$\ast\amalg\varnothing\to\ast\times\varnothing=\varnothing$
from a singleton to the empty set, which does not exist.}
Show that the identity canonically refines to a lax monoidal functor
$(\mathcal C,\times,\ast)\to(\mathcal C,\amalg,\varnothing)$.
When is this functor monoidal? Furthermore, show that any functor
$F\colon\mathcal C\to\mathcal D$ refines canonically to a lax symmetric monoidal functor
$(\mathcal C,\amalg,\varnothing)\to(\mathcal D,\amalg,\varnothing)$.
When is it monoidal?

\begin{proof}
Let $0$ be a zero object and write $0_{A,B}\colon A\to0\to B$
for the zero map. In the direction required by
Definition~1.2.31, the identity functor has structure maps
$\mu_{A,B}\colon A\amalg B\to A\times B$ determined by
\[
 \mu_{A,B}\iota_A=\langle\mathrm{id}_A,0_{A,B}\rangle,
 \qquad
 \mu_{A,B}\iota_B=\langle0_{B,A},\mathrm{id}_B\rangle,
\]
and unit $\varnothing\to\ast$.
Zero maps compose to zero, so $\mu$ is natural.
For associativity, identify the source with $A\amalg B\amalg D$
and the target with $A\times B\times D$ using the canonical
associators. Both composites restrict to
$\langle\mathrm{id},0,0\rangle$,
$\langle0,\mathrm{id},0\rangle$ and
$\langle0,0,\mathrm{id}\rangle$ on the three summands,
respectively, so they agree. The unit diagrams restrict to
the identity on the $A$-summand and agree on the
initial summand by uniqueness.

The unit map is already an isomorphism. Thus this functor is
monoidal exactly when every $\mu_{A,B}$ is an isomorphism,
equivalently when $\mathcal C$ has finite biproducts: the
coproduct $A\amalg B$, with projections
$[\mathrm{id}_A,0_{B,A}]$ and $[0_{A,B},\mathrm{id}_B]$,
must also be a product. This equivalence follows by comparing
that pair of projections with $A\times B$ via $\mu_{A,B}$.

For the second assertion, let $\mathcal D$ have finite
coproducts and let $\iota_A,\iota_B$ be the coproduct injections
in $\mathcal C$. Define
\[
 \begin{aligned}
  \mu^F_{A,B}&=[F\iota_A,F\iota_B]
   \colon F(A)\amalg F(B)\longrightarrow F(A\amalg B),\\
  u^F&\colon\varnothing_{\mathcal D}
                    \longrightarrow F(\varnothing_{\mathcal C}).
 \end{aligned}
\]
Here $u^F$ is the unique map from the initial object.
Naturality follows from functoriality. Associativity and
symmetry follow by checking each coproduct summand:
both routes are $F$ of the same
canonical injection. The unit diagrams are the identity
on $F(A)$ and agree on the initial summand by uniqueness.
Hence $F$ is canonically lax symmetric monoidal.
Its structure maps are invertible exactly when $F$ preserves
finite coproducts. Indeed, invertibility transports the
coproduct universal properties to their images under $F$;
conversely, preservation makes these maps canonical
isomorphisms between coproducts, including the initial object.
\end{proof}
\end{exercise}
